% Part: lambda-calculus
% Chapter: introduction
% Section: lambda-definability

\documentclass[../../../include/open-logic-section]{subfiles}

\begin{document}

\olfileid[tr]{lam}{int}{rep}
\olsection{Aritmetik Fonksiyonların \usetoken{S}{lambda definable} Olması}

Lambda kalkülüsü nasıl bir hesaplama modeli olabilir? Başta bu ifadeye nasıl
anlam verileceği bile açık değildir. Doğal sayılar üzerindeki
hesaplanabilirlikten söz etmek için bu sayılara uygun bir temsil bulmalıyız.
Şaşırtıcı derecede iyi çalışan bir temsil aşağıdadır.

\begin{defn}
Her $n$ doğal sayısı için \emph{Church sayısı} $\num{n}$, toplam $n$ tane
$x$ bulunan $\lambd[x][\lambd[y][(x(x(x(\dots x(y))))]]$ lambda terimi
olarak tanımlansın.
\end{defn}

$\num{n}$ terimleri ``yineleyicilerdir'': $f$ girdisinde $\num{n}$, $y$'yi
$f^n(y)$'ye götüren fonksiyonu döndürür. Her sayının normal olduğuna dikkat
edin. Artık bir lambda teriminin doğal sayılar üzerindeki bir fonksiyonu
``hesaplamasının'' ne demek olduğunu söyleyebiliriz.

\begin{defn}
$f(x_0,\dots,x_{k-1})$, $\Nat$'tan $\Nat$'a giden $n$-li kısmi bir fonksiyon
olsun. Bir $F$ $\lambd$-terimi her $n_0$, \dots, $n_{k-1}$ doğal sayı dizisi
için, $f(n_0,\dots,n_{k-1})$ tanımlı olduğunda
\[
F\, \num{n_0}\, \num{n_1} \dots \num{n_{k-1}} \red \num{f(n_0, n_1, \dots,
  n_{k-1})},
\]
tanımsız olduğunda ise $F\,\num{n_0}\,\num{n_1}\dots\num{n_{k-1}}$ normal
biçime sahip değilse, $F$'nin $f$'yi \emph{!!{lambda define}s} olduğunu
söyleriz.
\end{defn}

\begin{thm}
\ollabel{thm:lambda-def}
Bir $f$ fonksiyonu ancak ve ancak bir lambda terimi tarafından
!!{lambda defined} ise kısmi hesaplanabilir bir fonksiyondur.
\end{thm}

\begin{explain}
Bu teorem oldukça çarpıcıdır. Bir hesaplama modeli olarak lambda kalkülüsü
çok yalın bir kalkülüstür; yalnızca lambda soyutlaması ve uygulama işlemleri
vardır!{} Yine de bu kıt kaynaklarla her hesaplama yordamını gerçekleştirmek
mümkündür.
\end{explain}

\end{document}
